\documentclass[aspectratio=169,11pt]{beamer}

\usepackage[spanish,es-nodecimaldot]{babel}
\usepackage[T1]{fontenc}
\usepackage[utf8]{inputenc}
\usepackage{tikz}
\usetikzlibrary{arrows.meta,positioning}
\usepackage{xcolor}
\usepackage{amsmath}
\usepackage{booktabs}
\usepackage{listings}

\definecolor{azul}{HTML}{174A7E}
\definecolor{azulclaro}{HTML}{DCECF8}
\definecolor{verde}{HTML}{2E7D32}
\definecolor{verdeclaro}{HTML}{E8F5E9}
\definecolor{naranja}{HTML}{D95F02}
\definecolor{rojo}{HTML}{B33A3A}
\definecolor{fondo}{HTML}{F7F8FA}

\usetheme{Madrid}
\setbeamercolor{structure}{fg=azul}
\setbeamercolor{frametitle}{fg=azul}
\setbeamercolor{normal text}{fg=black!80}
\setbeamercolor{block title}{bg=azul,fg=white}
\setbeamercolor{block body}{bg=azulclaro,fg=black!80}
\setbeamertemplate{navigation symbols}{}
\setbeamertemplate{footline}{\hfill\insertframenumber\hspace{0.4cm}\vspace{0.2cm}}

\title{Laboratorio 8}
\subtitle{Redes con memoria sensorial para romper ciclos en navegacion reactiva}
\author{Inteligencia Artificial}
\date{}

\begin{document}

\begin{frame}
  \titlepage
\end{frame}

\begin{frame}{Motivacion}
  \begin{block}{Problema del robot reactivo}
    En el laboratorio anterior, el robot respondia a la observacion actual:
    \[
    a_t = \pi(s_t)
    \]
  \end{block}
  \begin{alertblock}{Limitacion}
    Si la secuencia de decisiones se repite, el robot puede quedar atrapado en un ciclo aunque la situacion local no haya cambiado.
  \end{alertblock}
\end{frame}

\begin{frame}{Del estado instantaneo al estado temporal}
  La idea clave es que la decision depende no solo del presente, sino tambien del pasado reciente:
  \[
  a_t = \pi(s_t, s_{t-1}, a_{t-1})
  \]
  donde:
  \begin{itemize}
    \item $s_t$: percepcion actual
    \item $s_{t-1}$: percepcion anterior
    \item $a_{t-1}$: accion ejecutada en el paso anterior
  \end{itemize}
  \vspace{0.3cm}
  \begin{block}{Resultado}
    La red ya no decide solo con el presente, sino con memoria de corto plazo.
  \end{block}
\end{frame}

\begin{frame}{El ciclo sensorimotor}
  \begin{center}
    \begin{tikzpicture}[node distance=1.7cm]
      \node[draw,rounded corners,fill=azulclaro] (entorno) {Entorno};
      \node[draw,rounded corners,fill=verde!15,right=of entorno] (sensores) {Sensores};
      \node[draw,rounded corners,fill=orange!20,right=of sensores] (red) {Red neuronal};
      \node[draw,rounded corners,fill=azulclaro,right=of red] (accion) {Accion};
      \draw[->,very thick,azul] (entorno) -- (sensores);
      \draw[->,very thick,azul] (sensores) -- (red);
      \draw[->,very thick,azul] (red) -- (accion);
      \draw[->,very thick,azul] (accion) to[bend left=35] (entorno);
    \end{tikzpicture}
  \end{center}
  \vspace{0.3cm}
  \begin{block}{Idea}
    El entorno genera sensaciones; el agente actua; la accion modifica el entorno; el ciclo se repite en $t, t+1, t+2, \dots$
  \end{block}
\end{frame}

\begin{frame}{Que es $t-1$ y por que importa}
  El robot necesita recordar informacion reciente:
  \[
  x_t = [s_t, s_{t-1}, a_{t-1}]
  \]
  \begin{itemize}
    \item $s_t$ representa el presente.
    \item $s_{t-1}$ aporta contexto temporal.
    \item $a_{t-1}$ informa que la red ya eligio antes.
  \end{itemize}
  \vspace{0.2cm}
  \begin{alertblock}{Consecuencia}
    La politica puede detectar secuencias repetitivas y evitar caer en ciclos locales.
  \end{alertblock}
\end{frame}

\begin{frame}{Red con memoria sensorial}
  \begin{block}{Arquitectura}
    \begin{center}
      12 entradas $\longrightarrow$ 16 neuronas ReLU $\longrightarrow$ 4 acciones
    \end{center}
  \end{block}
  \begin{itemize}
    \item \textbf{Entradas:} sensores actuales, sensores anteriores y accion previa.
    \item \textbf{Capa oculta:} combina la informacion para aprender patrones locales.
    \item \textbf{Salidas:} un valor por cada accion $\{U,D,L,R\}$.
    \item \textbf{Decision:} se elige la salida con mayor valor mediante $\arg\max$.
  \end{itemize}
  \begin{alertblock}{Interpretacion}
    La red usa el presente y el paso anterior para decidir hacia donde moverse.
  \end{alertblock}
\end{frame}

\begin{frame}{Formulacion matematica}
  La red calcula:
  \[
  h_t = \mathrm{ReLU}(W_1 x_t + b_1)
  \]
  \[
  y_t = W_2 h_t + b_2
  \]
  \[
  a_t = \arg\max(y_t)
  \]
  con:
  \[
  x_t = [s_t, s_{t-1}, a_{t-1}]
  \]
  \vspace{0.3cm}
  La accion elegida es la de mayor valor logit. La red compara todas las opciones con el estado presente y reciente.
\end{frame}

\begin{frame}{Por que ayuda a romper bucles}
  Un bucle aparece cuando la secuencia:
  \[
  s_t, a_t, s_{t+1}, a_{t+1}, \dots
  \]
  vuelve a un estado similar ya observado.
  \begin{block}{Con memoria}
    El agente puede reconocer que:
    \begin{itemize}
      \item la observacion actual es parecida a la anterior,
      \item la accion anterior fue la misma,
      \item la secuencia se vuelve circular.
    \end{itemize}
  \end{block}
\end{frame}

\begin{frame}{Como se evalua}
  La estrategia sigue siendo evolutiva:
  \begin{enumerate}
    \item generar 40 individuos,
    \item evaluar fitness,
    \item conservar elites,
    \item cruzar y mutar,
    \item seleccionar la mejor red.
  \end{enumerate}
  \vspace{0.2cm}
  \begin{block}{Fitness}
    \[
    F = 500 - 35d + 4p - 30c - 3r + 2000m - 10b
    \]
    donde $b$ es una penalizacion por bucles.
  \end{block}
\end{frame}

\begin{frame}{Hipotesis de trabajo}
  \begin{block}{Hipotesis}
    Una red con entrada de memoria sensorial y accion previa reducira la frecuencia de ciclos y mejorara la capacidad de navegacion local en entornos con obstaculos.
  \end{block}
  \vspace{0.25cm}
  \begin{center}
    \textbf{Objetivo experimental: comparar la red basica contra la red con memoria.}
  \end{center}
\end{frame}

\begin{frame}{Experimentos sugeridos}
  \begin{itemize}
    \item Red basica: 4 sensores actuales.
    \item Red con memoria: 4 actuales + 4 previos + accion previa.
    \item Mapa sencillo, mapa con pasillos estrechos y mapa con bucles locales.
  \end{itemize}
  \vspace{0.3cm}
  \begin{block}{Metricas}
    exito, choques, pasos, ciclos, distancia final.
  \end{block}
\end{frame}

\begin{frame}{Conclusiones}
  \begin{itemize}
    \item La memoria sensorial hace que la decision sea menos reactiva pura.
    \item El robot valora no solo lo que ve, sino lo que ya hizo.
    \item La red aprende una politica con dependencia temporal.
    \item Esto es una forma simple de inteligencia sensoriomotora.
  \end{itemize}
\end{frame}

\end{document}
